\chapter{The Six Operations with Complete Coefficients}
\label{chap:proetale-six}

Let \(f:X\to Y\) be a separated morphism of finite type between noetherian
schemes, and let \(\Lambda=\Zl\) or \(\mathcal O_E\), with \(\ell\) invertible
on both schemes.  We now construct the six operations on
\(\Dcons(-,\Lambda)\).  The construction separates two logically different
ingredients:
\begin{itemize}
\item the pro-etale completion formalism, which passes coherently from
      torsion levels to one integral category;
\item the deep torsion theorems---finiteness, base change, compactification,
      and purity---which the change of site does not prove.
\end{itemize}

\section{Completed pullback and direct image}

\begin{definition}[Completed pullback]
\label{def:completed-pullback}
For a complete complex \(K\) on \(Y\), define
\begin{equation}
\label{eq:completed-pullback}
 \Lhat f^*K
   :=\widehat{\mathbf L f^*K}
   \simeq
   \Rlim_n \mathbf L f^*(K\otimes^{\mathbf L}_{\Lambda}\Lambda_n).
\end{equation}
The hat is part of the operation: the naive inverse-image complex need not be
derived complete.
\end{definition}

\begin{proposition}[Adjunction and constructibility]
\label{prop:completed-pullback-adjunction}
The functor \(\Lhat f^*\) preserves lisse and constructible complexes and is
left adjoint to \(\R f_*\) on complete complexes:
\[
 \operatorname{Map}_X(\Lhat f^*K,L)
   \simeq \operatorname{Map}_Y(K,\R f_*L).
\]
\end{proposition}

\begin{proof}
The usual inverse image is left adjoint to direct image in the ambient
ringed topoi.  Applying the completion reflector does not change maps into a
complete target, which proves the adjunction.  Pullback of a finite locally
constant complex is finite locally constant; a common finite stratification
therefore proves preservation of constructibility.  This is
\cite[Lemma~6.5.9]{BhattScholzeProEtale}.
\end{proof}

\begin{warning}[Why the completion sign cannot be suppressed]
For special morphisms and bounded objects, ordinary pullback may already be
complete.  It is not true uniformly.  Defining the integral operation as
ordinary pullback and later asserting completeness would make functoriality
depend on extra hypotheses.  The canonical operation is
\(\Lhat f^*\), as emphasized in
\cite[Remark~6.5.10]{BhattScholzeProEtale}.
\end{warning}

The hard input for direct images is first a finite-coefficient theorem.

\begin{theorem}[Torsion finiteness input]
\label{thm:torsion-finiteness-input}
Let \(A\) be a finite \(\ell\)-primary ring.  For a separated finite-type
morphism \(f:X\to Y\),
\[
 \R f_*:\D^b_c(X,A)\longrightarrow\D^b_c(Y,A)
\]
preserves constructibility and has a cohomological-amplitude bound depending
on \(f\) but not on the exponent annihilating \(A\).
\end{theorem}

\begin{proof}
This is the Grothendieck--Gabber finiteness theorem.  Its proof uses
compactification, cohomological-dimension bounds for affine schemes, proper
base change, and noetherian induction; it is not a formal consequence of the
pro-etale topology.  The precise package used in the passage to complete
coefficients is \cite[Theorem~6.7.1]{BhattScholzeProEtale}, based on the
torsion theory assembled in \cite{SGA4,SGA4half,ILOTravauxGabber}.
\end{proof}

\begin{corollary}[Integral direct image]
\label{cor:integral-direct-image}
Direct image restricts to
\[
 \R f_*:\Dcons(X,\Lambda)\longrightarrow\Dcons(Y,\Lambda),
\]
and reduction commutes with it:
\[
 ( \R f_*K)\otimes^{\mathbf L}_{\Lambda}\Lambda_n
   \simeq
 \R f_*(K\otimes^{\mathbf L}_{\Lambda}\Lambda_n).
\]
\end{corollary}

\begin{proof}
Apply the torsion theorem to every reduction \(K_n\).  Its uniform amplitude
bound permits \(\R f_*\) to commute with the homotopy limit of the compatible
tower.  The resulting limit is complete by
\cref{thm:proetale-complete-system-equivalence}, and its first reduction is
constructible.  This is \cite[Lemma~6.7.2]{BhattScholzeProEtale}.
\end{proof}

\section{Base change}

Consider a cartesian square
\begin{equation}
\label{eq:six-base-change-square}
\begin{tikzcd}
X' \arrow[r,"g'"] \arrow[d,"f'"'] & X \arrow[d,"f"] \\
Y' \arrow[r,"g"'] & Y .
\end{tikzcd}
\end{equation}

\begin{theorem}[Integral base-change theorems]
\label{thm:proetale-base-change}
For constructible complete complexes:
\begin{enumerate}
\item if \(g\) is smooth, then
\[
 \Lhat g^*\R f_*K
   \xrightarrow{\sim}
 \R f'_*\Lhat {g'}^*K;
\]
\item if \(f\) is proper, the same map is an equivalence for arbitrary
      \(g\);
\item for \(\R f_!\), constructed below, arbitrary base change holds:
\[
 \Lhat g^*\R f_!K
   \xrightarrow{\sim}
 \R f'_!\Lhat {g'}^*K.
\]
\end{enumerate}
\end{theorem}

\begin{proof}
Reduce modulo \(\mathfrak m^n\), apply respectively smooth base change,
proper base change, and compact-support base change for finite torsion
coefficients, and pass to the inverse limit.  Uniform cohomological bounds
justify the passage.  The integral statements are
\cite[Lemmas~6.7.4, 6.7.5, and 6.7.10]{BhattScholzeProEtale}.
\end{proof}

\section{Compact support is a derived geometric construction}

By Nagata compactification, factor \(f\) as
\[
 X\xrightarrow{j}\overline X\xrightarrow{\bar f}Y,
\]
where \(j\) is an open immersion and \(\bar f\) is proper.

\begin{definition}[Exceptional direct image]
\label{def:proetale-fshriek}
Define
\begin{equation}
\label{eq:fshriek-compactification}
 \R f_!K:=\R\bar f_*\,j_!K.
\end{equation}
Here extension by zero \(j_!\) is taken levelwise on torsion coefficients and
then completed.
\end{definition}

\begin{theorem}[Independence of compactification]
\label{thm:fshriek-independent}
Definition~\eqref{eq:fshriek-compactification} is independent, up to a
unique coherent equivalence, of the compactification.  It is functorial in
\(f\), preserves constructible complexes, equals \(\R f_*\) for proper
\(f\), and satisfies the arbitrary base-change formula in
\cref{thm:proetale-base-change}.
\end{theorem}

\begin{proof}
Any two compactifications are dominated by the schematic closure of the
diagonal image of \(X\) in their fiber product.  Proper base change and the
identity \(p_*j'_!\simeq j_!\) for the resulting proper modification compare
the two definitions.  Passing to a third dominating compactification proves
the cocycle condition.  This is
\cite[Definition~6.7.6 and Lemmas~6.7.7 and 6.7.10]
{BhattScholzeProEtale}.
\end{proof}

\begin{warning}[\(f_!\) is not obtained by deriving an abelian functor]
For a general separated finite-type morphism, the operation in
\cref{def:proetale-fshriek} is defined in the derived category through a
compactification.  It is not the left or right derived functor of a globally
defined abelian-level operation \(f_!\) having the required properties.
Only in special cases, such as an open immersion, is there an exact
abelian-level extension-by-zero functor.  Treating general \(\R f_!\) as
\(\R\) applied to an underived functor loses the construction and its
coherence; see \cite[Remarks~6.7.8--6.7.9]{BhattScholzeProEtale}.
\end{warning}

\begin{proetalebox}
The pro-etale simplification is now visible: the compactification argument is
performed at finite torsion level only once, and repleteness turns the
compatible family into one functor on complete \(\Zl\)-complexes.  It does
not remove Nagata compactification or proper base change.
\end{proetalebox}


\section{Tensor product, internal Hom, and projection formulas}

\begin{proposition}
\label{prop:proetale-tensor-hom}
For \(K,L\in\Dcons(X,\Lambda)\), the objects
\[
 K\widehat\otimes^{\mathbf L}_{\Lambda}L,
 \qquad
 \R\underline{\operatorname{Hom}}_{\Lambda}(K,L)
\]
are constructible.  They form a closed monoidal pair, and for
\(M\in\Dcons(Y,\Lambda)\) there is a projection-formula equivalence
\begin{equation}
\label{eq:fshriek-projection}
 \R f_!(K\widehat\otimes^{\mathbf L}\Lhat f^*M)
   \simeq
 (\R f_!K)\widehat\otimes^{\mathbf L}M.
\end{equation}
\end{proposition}

\begin{proof}
On each stratum this is perfect-complex algebra, so tensor and internal Hom
remain lisse.  The projection formula is checked modulo
\(\mathfrak m^n\), where it is the torsion projection formula, and then
completed.  See \cite[Lemmas~6.7.12--6.7.14]{BhattScholzeProEtale}.
\end{proof}

\begin{warning}
The analogous formula with \(\R f_*\) in place of \(\R f_!\) is not valid
without additional hypotheses; properness is a sufficient one.  This is
\cite[Remark~6.7.15]{BhattScholzeProEtale}.
\end{warning}

\section{Exceptional pullback and Verdier duality}

\begin{theorem}[Existence of \(f^!\)]
\label{thm:proetale-fupper-shriek}
The functor \(\R f_!\) has a right adjoint
\[
 f^!:\Dcons(Y,\Lambda)\longrightarrow\Dcons(X,\Lambda).
\]
Let \(S\) carry a dualizing object \(\omega_S\), and suppose \(X\) and
\(Y\) are separated of finite type over \(S\), with structural maps
\(a_X\) and \(a_Y\).  Set
\[
 \omega_Z=a_Z^!\omega_S,
 \qquad
 \mathbb D_ZK=
 \R\underline{\operatorname{Hom}}_{\Lambda}(K,\omega_Z)
 \quad(Z\in\{X,Y\}).
\]
Then the biduality morphism
\[
 K\longrightarrow\mathbb D_X\mathbb D_XK
\]
is an equivalence for constructible \(K\), and
\begin{equation}
\label{eq:six-duality-identities}
 f^!\simeq\mathbb D_X\Lhat f^*\mathbb D_Y,
 \qquad
 \R f_!\simeq\mathbb D_Y\R f_*\mathbb D_X.
\end{equation}
\end{theorem}

\begin{proof}
At torsion level, duality is constructed from a compactification, proper
duality, and the exceptional inverse image for a closed immersion.  Uniform
amplitude bounds make the inverse limit complete and preserve the adjunction.
Biduality is local for a finite stratification and reduces on lisse objects
to biduality for perfect modules.  The integral construction is
\cite[Lemmas~6.7.11 and 6.7.19--6.7.20]{BhattScholzeProEtale}.
\end{proof}

For a closed immersion \(i:Z\hookrightarrow X\) with open complement
\(j:U\hookrightarrow X\), the adjunctions give the two localization triangles
\begin{align*}
 j_!j^*K&\longrightarrow K\longrightarrow i_*i^*K\xrightarrow{+1}, \\
 i_*i^!K&\longrightarrow K\longrightarrow\R j_*j^*K\xrightarrow{+1}.
\end{align*}
These triangles, rather than stalkwise case divisions, will later drive
induction on supports.

\begin{theorem}[Purity inputs]
\label{thm:proetale-purity-inputs}
Let coefficients be as above.
\begin{enumerate}
\item If \(f:X\to Y\) is smooth of pure relative dimension \(d\), then
\[
 f^!K\simeq\Lhat f^*K(d)[2d].
\]
\item If \(i:Z\hookrightarrow X\) is a regular immersion of pure codimension
\(c\) between regular noetherian schemes, then
\[
 i^!\Lambda\simeq\Lambda(-c)[-2c].
\]
\end{enumerate}
\end{theorem}

\begin{proof}
The first assertion is relative cohomological purity and the second is
Gabber's absolute purity theorem.  Both are proved for torsion coefficients
using deformation to the normal cone and the calculation for the zero
section of affine space; compatibility in \(n\) then gives the complete
coefficient statement.  These are deep geometric inputs, not consequences
of completion.  For the coherent six-functor formulation see
\cite[Theorems~7.15, 7.16, and 7.19]{ScholzeSixFunctors}; for absolute purity
see \cite{ILOTravauxGabber}.
\end{proof}

\section{The coherent package}

\begin{theorem}[Constructible pro-etale six operations]
\label{thm:proetale-six-functor-package}
For noetherian schemes with \(\ell\) invertible and separated finite-type
morphisms, the assignment
\[
 X\longmapsto\Dcons(X,\Lambda)
\]
carries the operations
\[
 \Lhat f^*,\quad \R f_*,\quad \R f_!,\quad f^!,
 \quad -\widehat\otimes^{\mathbf L}-,
 \quad \R\underline{\operatorname{Hom}}(-,-)
\]
with their adjunctions, composition equivalences, base-change maps,
projection formulas, localization triangles, and Verdier duality.  All
coherence diagrams commute.  Scalar extension along
\(\Zl\to\mathcal O_E\to E\) commutes with the six operations.
\end{theorem}

\begin{proof}
Existence and the individual equivalences were proved above by compatible
finite-level constructions.  Coherence can also be checked modulo every
\(\mathfrak m^n\): reduction is conservative on complete objects, and the
torsion six-functor formalism is coherent.  Rational coefficients follow by
localization.  An infinity-categorical construction that packages all
coherences at once is \cite[Theorem~7.15]{ScholzeSixFunctors}.
\end{proof}

\begin{foundation}
The logical dependency of the theorem is
\[
\begin{tikzcd}[column sep=normal,row sep=small]
\substack{\text{torsion finiteness}\\
          \text{and a uniform dimension bound}}
  \arrow[r] & \text{integral }\R f_* \arrow[d] \\
\substack{\text{Nagata compactification}\\
          \text{and proper base change}}
  \arrow[r] & \text{integral }\R f_! \arrow[d] \\
\text{torsion duality and purity}
  \arrow[r] & \substack{\text{integral }f^!\\
                         \text{and Verdier duality}}.
\end{tikzcd}
\]
Repleteness and derived completion justify every vertical passage from
finite coefficients to complete coefficients.  They do not supply any of
the three inputs in the left column.
\end{foundation}

\begin{formalization}
The six operations should be encoded as a coherent functorial package, but
its theorems should be exposed through small interfaces: adjunction,
base-change square, projection formula, and localization triangle.  The
definition of \(\R f_!\) must remember the proof that compactification choices
produce canonically equivalent functors.  Completed pullback and ordinary
pullback should never be definitionally identified.
\end{formalization}

\begin{auditbox}
The pro-etale site eliminates the adic quotient category and makes all six
operations act on honest complete objects.  It does not simplify away the
Grothendieck--Gabber finiteness theorem, Nagata compactification, proper and
smooth base change, Verdier duality, or absolute purity.  A proof that lists
only the six operations on the pro-etale site without these inputs has hidden
the hardest foundational work.
\end{auditbox}
